\documentclass[runningheads,a4paper]{llncs}

\usepackage[T1]{fontenc}
\usepackage{lmodern}
\usepackage{fullpage}
\usepackage{microtype}
\usepackage{amsmath,amssymb,mathtools}
\usepackage[hidelinks]{hyperref}
\usepackage[nameinlink,noabbrev]{cleveref}

% Shared notation belongs here once it has been fixed and reviewed.
\newcommand{\bits}{\{0,1\}}
\newcommand{\negl}{\mathop{\mathsf{negl}}}
\newcommand{\poly}{\mathop{\mathsf{poly}}}
\newcommand{\GC}{\mathsf{GC}}
\newcommand{\pk}{\mathsf{pk}}
\newcommand{\sk}{\mathsf{sk}}
\newcommand{\sample}{\mathrel{\leftarrow}}

% Algorithm names, game commands, and record tags.
\newcommand{\Setup}{\mathsf{Setup}}
\newcommand{\Sample}{\mathsf{Sample}}
\newcommand{\IdealSample}{\mathsf{IdealSample}}
\newcommand{\QueryRO}{\mathsf{QueryRO}}
\newcommand{\Declare}{\mathsf{Declare}}
\newcommand{\DeclareOutput}{\mathsf{DeclareOutput}}
\newcommand{\Finish}{\mathsf{Finish}}
\newcommand{\SimSetup}{\mathsf{SimSetup}}
\newcommand{\SimRO}{\mathsf{SimRO}}
\newcommand{\SimSample}{\mathsf{SimSample}}
\newcommand{\Real}{\mathsf{Real}}
\newcommand{\Ideal}{\mathsf{Ideal}}
\newcommand{\Inconsistent}{\mathsf{inconsistent}}
\newcommand{\SimState}{\mathsf{st}}

\newcommand{\Garble}{\mathsf{Garble}}
\newcommand{\Encode}{\mathsf{Encode}}
\newcommand{\Eval}{\mathsf{Eval}}
\newcommand{\Topo}{\mathsf{Topo}}
\newcommand{\GarbleAux}{\mathsf{aux}}
\newcommand{\SimGarble}{\mathsf{SimGarble}}
\newcommand{\SimEncode}{\mathsf{SimEncode}}
\newcommand{\GarbleReal}{\mathsf{GarbleReal}}
\newcommand{\GarbleIdeal}{\mathsf{GarbleIdeal}}
\newcommand{\GarbleCircuit}{\mathsf{GarbleCircuit}}
\newcommand{\EncodeInput}{\mathsf{EncodeInput}}
\newcommand{\PrvOne}{\mathsf{prv1}}

\newcommand{\KeyGen}{\mathsf{KeyGen}}
\newcommand{\Enc}{\mathsf{Enc}}
\newcommand{\Dec}{\mathsf{Dec}}
\newcommand{\SimNCE}{\mathsf{SimNCE}}
\newcommand{\ExplainNCE}{\mathsf{ExplainNCE}}
\newcommand{\GetPublicKey}{\mathsf{GetPublicKey}}
\newcommand{\EncryptMessage}{\mathsf{EncryptMessage}}
\newcommand{\CorruptReceiver}{\mathsf{CorruptReceiver}}
\newcommand{\CorruptSender}{\mathsf{CorruptSender}}
\newcommand{\NCEReal}{\mathsf{NCEReal}}
\newcommand{\NCEIdeal}{\mathsf{NCEIdeal}}

\newcommand{\Share}{\mathsf{Share}}
\newcommand{\Reconstruct}{\mathsf{Reconstruct}}
\newcommand{\ShareSecret}{\mathsf{ShareSecret}}
\newcommand{\ReadShare}{\mathsf{ReadShare}}
\newcommand{\ShareGame}{\mathsf{ShareGame}}
\newcommand{\CompleteShares}{\mathsf{CompleteShares}}

\newcommand{\CPAGame}{\mathsf{CPAGame}}
\newcommand{\Challenge}{\mathsf{Challenge}}

% Ledger notation.
\newcommand{\FLedger}{\mathcal F_{\mathsf{Ledger}}}
\newcommand{\crs}{\mathsf{crs}}
\newcommand{\WhosTurn}{\mathsf{WhosTurn}}
\newcommand{\NoMsg}{\mathsf{NoMsg}}

% Application interface.
\newcommand{\FMPC}{\mathcal F_{\mathsf{MPC}}}
\newcommand{\pid}{\mathsf{pid}}
\newcommand{\ncek}{\mathsf{ncek}}

% State-transfer keys, distinct from registered delivery keys.
\newcommand{\statepk}{\mathsf{pk}^{\mathsf{state}}}
\newcommand{\statesk}{\mathsf{sk}^{\mathsf{state}}}

% Sample-owned input-encryption keys and registered signing keys.
\newcommand{\inputpk}{\mathsf{pk}^{\mathsf{in}}}
\newcommand{\inputsk}{\mathsf{sk}^{\mathsf{in}}}
\newcommand{\sigvk}{\mathsf{vk}^{\mathsf{sig}}}
\newcommand{\sigsk}{\mathsf{sk}^{\mathsf{sig}}}
\newcommand{\sig}{\mathsf{sig}}
\newcommand{\sid}{\mathsf{sid}}
\newcommand{\SigGen}{\mathsf{SigGen}}
\newcommand{\Sign}{\mathsf{Sign}}
\newcommand{\Verify}{\mathsf{Verify}}
\newcommand{\GetVerificationKey}{\mathsf{GetVerificationKey}}
\newcommand{\SignMessage}{\mathsf{SignMessage}}
\newcommand{\Forge}{\mathsf{Forge}}
\newcommand{\InputRecord}{\mathsf{InputRecord}}
\newcommand{\DefaultInput}{\mathsf{DefaultInput}}


\title{The YAMSO Model: YOSO without Erasure}
\titlerunning{The YAMSO Model}
\author{Jesper Buus Nielsen\inst{1} \and Ivan Damg{\aa}rd\inst{1} \and
  Sannidhi V. Hebbar\inst{2}}
\authorrunning{J. B. Nielsen et al.}
\institute{Aarhus University, Denmark\\
  \email{\{jbn,ivan\}@cs.au.dk}
  \and
  Indian Institute of Science, India\\
  \email{sannidhiv@iisc.ac.in}}

\begin{document}

\maketitle

\begin{abstract}
The YOSO paradigm protects against a strongly adaptive adversary. Each
ephemeral role prepares one message, erases all remaining state, and only
then speaks. We ask what remains possible when erasure is unavailable and
speaking instead exposes the sender's complete state. At first sight, such a
model appears useless. Nevertheless, it has a different irreversible
resource: a machine may choose never to speak, in which case its state and
even its role can remain hidden. We call the resulting model YAMSO, for ``You
At Most Speak Once.'' This preliminary article introduces the model,
describes a basic non-reactive construction from anonymous hidden committees
and garbled circuits, and isolates the additional correlated setup needed for
ongoing computation.
\keywords{adaptive security \and YOSO \and erasure-free protocols \and
anonymous committees \and garbled circuits}
\end{abstract}

\section{Introduction}
\label{sec:introduction}

The YOSO (You Only Speak Once) model was introduced as a way to obtain secure
multiparty computation against a very strong adaptive adversary
\cite{GentryEtAl2021YOSOFull}. A protocol role is assigned to an ephemeral
party which prepares one message, deletes the secret state used to prepare
it, and then sends the message. Only at this point does the role reveal
itself. Corrupting the sender after publication is therefore too late: the
message is already available and the secrets that produced it no longer
exist. Later work has developed this anonymous, ephemeral-role paradigm in
several directions; the erase-before-send discipline remains central.

This paper starts from the deliberately unfriendly question: what if deletion
is impossible? The restriction might be legal---a regime may require state
to be retained for later inspection---or technological. For example, a
party running in a cloud environment might fear that speaking identifies its
machine and enables the cloud operator to inspect a residual snapshot despite
an attempted deletion. We capture the extreme version of this concern by
declaring that whenever an ordinary party speaks, its complete state leaks to
the adversary.

The immediate answer seems to be ``obviously impossible.'' Eventually all
parties relevant to the computation will speak and hence reveal everything.
This conclusion overlooks one asymmetry. Speaking is optional. A machine
may remain silent forever, in which case its secret state may stay hidden. We
therefore place mutually exclusive pieces of state behind different
public-key roles. The selected roles speak and are expended; the unselected
roles never speak. In this sense, silence replaces erasure as the one-time
resource.

We call the resulting setting \emph{YAMSO}: You At Most Speak Once. In a
YAMSO execution, a synchronous ledger maintains a growing ordered list of
public keys and messages. The public-key index determines the protocol role.
For the preliminary model, each newly registered computation role is corrupted
independently with probability $p$, whereas named application parties remain
adaptively corruptible in the usual way. An honest role keeps its secret key
inside a protected ledger interface until its one scheduled opportunity to
speak.
The intended later model will derive essentially the same random committee
corruption from hidden assignment against an adversary controlling a bounded
fraction of a much larger population. Once a computation role speaks, its key
may become completely public; its message is already irrevocably recorded and
the role is never used again.

This first version has three modest goals. First, it gives a generic ledger
API with synchronous public state, growing registration, random corruption,
protected local evaluation, and explicit one-time turns. Second, it gives a
finite construction sketch that illustrates why the model is not vacuous:
each alternative garbled input label is shared behind a different committee,
and only the committee selected by the input token reveals its keys. Third, it
proposes a \emph{universal streamer}. This changes the universal-sampler
interface by giving a fixed sampling description $d$ an append-only stream of
public inputs. The outputs for successive input prefixes must themselves be
prefixes. A fixed $d$ can thereby describe one indefinitely extended garbling
whose output state labels are literally the input state labels of the next
stretch, while later public keys arrive as public inputs. Constructing such a
prefix-preserving universal sampler remains open.

The presentation is intentionally preliminary. The construction in
\cref{sec:basic-construction} is a feasibility blueprint rather than a stated
theorem. Section~\ref{sec:model} fixes an operational experiment against
which a first proof can be written. The detailed UC functionality is deferred
until this API is stable, and universal sampling remains entirely a
protocol-level mechanism.

\paragraph{Relation to prior work.}
Our starting point is YOSO \cite{GentryEtAl2021YOSOFull}. Anonymous
blockchain committees tolerating a corruption rate around one quarter are
known, but their secret handoff relies on erasures
\cite{BenhamoudaEtAl2020BlockchainSecret}. Blockchain one-time programs
already demonstrate conditional release of one garbled input on an immutable
ledger \cite{GoyalGoyal2017Blockchains}. No-erasure hidden committees also
appear in communication-local MPC, although there the selected committee runs
an ordinary multi-round computation
\cite{ChandranEtAl2024CommunicationLocalMPC}. Our basic sketch combines the
one-time-role viewpoint with threshold-separated garbled labels. The
universal-streamer direction connects to reactive garbling, succinct
garbling, garbled Turing machines, and streaming functional encryption
\cite{NielsenRanellucci2015ReactiveGarbling,LinPass2014SuccinctGarbling,%
GoldwasserEtAl2013TuringMachines,GuanKorbSahai2022StreamingFE}.

\section{Preliminaries: What Must Be Fixed}
\label{sec:preliminaries}

This section is presently a roadmap. The final paper should reproduce every
definition whose details affect the YAMSO claim, while importing routine
cryptographic syntax and standard security definitions by reference.

\subsection{Conventions}

We write $\lambda$ for the security parameter, $[n]=\{1,\ldots,n\}$, and
$\bits^*$ for the set of finite bit strings. All adversaries and protocols
are probabilistic polynomial time unless explicitly stated otherwise, and
$\negl(\lambda)$ denotes a negligible function. Protocol identifiers,
epochs, wire identifiers, and ledger slots must be included in every
cryptographic domain separator.

\subsection{Definitions that belong in this paper}

The following objects are model-specific and should eventually be defined in
full.

\begin{description}
  \item[Ledger.] An append-only, totally ordered bulletin board with an
    explicit acceptance predicate. Each logical slot accepts at most one
    authenticated value; later protocol steps depend on the accepted prefix,
    not merely on the multiset of posted strings.

  \item[Anonymous machine registry.] A finite eligible population for each
    security experiment. Its public records are computationally
    indistinguishable with respect to hidden role assignment. A formulation
    may use independently registered public keys, or an anonymous IBE-style
    master public key with unlinkable, eligibility-limited secret-key
    extraction.

  \item[YAMSO machine.] A machine either appends one protocol message during
    its lifetime or remains silent forever. After an append, its complete
    prior state is exposed. Role reuse must be forbidden, since exposing one
    role would otherwise expose a future role stored by the same machine.

  \item[Strong input provider.] An exceptional party may append a prescribed
    function of its state rather than the state itself. If it is later
    corrupted, the ideal process releases its input to the simulator; if it
    is never corrupted, its input must remain private.

  \item[Role-oblivious corruption.] The adversary may adaptively inspect a
    bounded fraction of the machine universe, but unopened machines remain
    exchangeable. Discovering one committee member must not provide a test
    for identifying another unopened member.
\end{description}

\subsection{Primitives that can largely be imported}

Standard definitions of pseudorandom functions, commitments, signatures,
secret sharing, and circuit garbling need not be repeated in full. We should
state the exact variants used and cite them. The following nonstandard
interfaces do require care.

\begin{description}
  \item[Anonymous delivery.] Message privacy alone is insufficient: a ledger
    ciphertext must hide which registered key can decrypt it. We need a
    multi-user, key-private or anonymous encryption definition, together with
    correctness under recipient discovery. If ciphertexts are intended to
    look like unrelated ledger traffic, pseudorandom-ciphertext security must
    also be stated.

  \item[Verifiable threshold sharing.] The construction needs privacy below
    the reconstruction threshold, liveness from the honest members, and
    public rejection of malformed shares. We can import a standard scheme
    once the corruption threshold and synchrony assumptions are fixed.

  \item[Garbled circuits.] We can import the usual garbling syntax and privacy
    guarantee. The paper must separately prove that post-speech exposure of
    a selected label committee reveals no information beyond the label that
    is already reconstructed on the ledger.

  \item[Universal sampling.] The setup may be phrased using the universal
    sampler of Hofheinz et al.\ \cite{HofheinzEtAl2014UniversalSamplers}.
    Its ordinary samples for distinct descriptions use independent ideal
    randomness. Any ongoing construction that correlates different epochs
    must therefore expose that correlation explicitly rather than attribute
    it to the standard sampler definition.
\end{description}

\paragraph{Open drafting decision.}
For the first formal model it is probably cleanest to treat the ledger,
anonymous registry, and role assignment as ideal resources, while using
standard cryptographic primitives for label sharing and garbling. A later
section can then separate implementations of the setup---ordinary anonymous
public keys, anonymous IBE, or universal sampling---from the core YAMSO
phenomenon.

\section{The YAMSO Ledger Model}
\label{sec:model}

We describe the ledger interface at the level needed for the finite
construction in \cref{sec:basic-construction}. Our approach follows the
modelling style of \cite[Section~2.1]{DinsdaleYoungEtAl2020Afgjort}: we state
an intentionally strong and simple abstraction, isolate the properties used
by the protocol, and postpone its detailed realization. A later version will
express this interface as an ideal functionality
$\mathcal F_{\mathsf{Ledger}}$, in the style of the total-order-broadcast
functionality of \cite{DamgardEtAl2025AsyncYOSO}. There is no reason to fix
that UC bookkeeping before the API itself has stabilized.

The ledger is generic. It knows nothing about universal sampling, garbled
circuits, committees, input wires, or the meaning of protocol roles. Those
belong to the protocol using the ledger.

\subsection{Public state and initialization}

Fix a security parameter $\lambda$, a session identifier $\mathsf{sid}$, a
constant $p<1/3$, a generic setup distribution $\mathsf{Setup}$, a
key-generation algorithm $\mathsf{KeyGen}$, and a description of a deterministic
protocol-dependent algorithm $\mathsf{WhosTurn}$. The description of
$\mathsf{WhosTurn}$ is supplied to the ledger as a fixed initialization value;
it cannot be changed in response to the execution. The ledger samples
\[
  \mathsf{crs}\leftarrow\mathsf{Setup}(1^\lambda)
\]
and initializes the public append-only list as
$L=((\mathsf{Setup},\mathsf{crs}))$. The concrete protocol gives meaning to
$\mathsf{crs}$; the ledger does not interpret it.

We assume synchronous computation and an unrealistically strong broadcast
abstraction. After an entry is appended, every party sees exactly the same
list $L$ at the same time. An append is atomic, cannot be suppressed after
it occurs, and fixes the position of that entry. There are no delayed,
private, or party-local ledger views. This deliberately idealized assumption
removes consensus and network scheduling from the present question.

Apart from its own setup, registration, and timeout records, the ledger does
not enforce protocol syntax. In particular, an arbitrary message supplied by
a corrupt scheduled party is appended verbatim. The protocol reading $L$
decides which entries are well-formed and ignores malformed ones.

\subsection{Two registration paths}

The registry grows during the execution. There are two ways to add a key;
neither command takes a protocol role as input. The interface imposes no
a-priori bound on the registry, although any efficient execution creates only
polynomially many keys and every finite protocol instance uses a finite
prefix.

\paragraph{Computation-party registration.}
On a $\mathsf{Register}$ request from the adversary, the ledger samples
\[
  (pk_i,sk_i)\leftarrow\mathsf{KeyGen}(1^\lambda),
\]
where $i$ is the next computation-key index, creates an ITM whose PID is
$pk_i$, and appends $(\mathsf{Register},pk_i)$ to $L$. Independently it
samples $c_i\leftarrow\mathsf{Bernoulli}(p)$. If $c_i=1$, the new
computation party is corrupt and $sk_i$ is given to the adversary. Otherwise
it is honest, the ledger retains $sk_i$ as protected state, and the adversary
cannot corrupt this party later. Registration is indivisible: after
requesting it, the adversary cannot discard the public key because it
dislikes the corruption outcome.

\paragraph{Application-party registration.}
A fixed application party with ordinary UC identity $P$ may also input
$\mathsf{Register}$. The ledger samples
$(pk_P,sk_P)\leftarrow\mathsf{KeyGen}(1^\lambda)$, stores the association
$P\leftrightarrow(pk_P,sk_P)$, and appends
$(\mathsf{Register},P,pk_P)$ to $L$. It does not return $sk_P$ to an honest
$P$. There is no Bernoulli corruption experiment for such a party: the
adversary may corrupt $P$ adaptively in the ordinary UC sense. If $P$ is
already corrupt, or becomes corrupt later, the ledger releases $sk_P$ and
its retained record of $P$'s protected calls to the adversary. The target
ideal functionality supplies the simulator with any input or output leakage
prescribed upon that corruption.

Each PID registers at most once per session. Application parties are the
meaningful principals at the interface of the computation---for example, an
account owner supplying an input. Their identity need not and generally
cannot be hidden. Computation parties, by contrast, are protocol machinery.

A public protocol-level rule maps registration positions to roles. Thus the
seventh computation key may represent the eighth share-release role. This
mapping is not an argument to $\mathsf{Register}$ and is not interpreted by
the ledger.

\subsection{Turns and protected local evaluation}

At each synchronous step, the fixed algorithm $\mathsf{WhosTurn}(L)$ returns
either a registered public key that has not previously had a turn, or
$\bot$. It may schedule a computation key $pk_i$ or the registered key
$pk_P$ of an application party. A well-formed protocol chooses
$\mathsf{WhosTurn}$ so that it waits for enough registrations and schedules
each required role in the appropriate order. The ledger does not interpret
why a particular key is scheduled.

When a key $pk$ is scheduled, its owner has one synchronous opportunity to
speak. Write $L_{mathsf{pre}}$ for the list at the start of that turn. An
honest owner privately gives the ledger a finite polynomial-time algorithm
$A$. The ledger samples and retains any coins required by $A$, and evaluates
\[
  m\leftarrow A(sk,L_{\mathsf{pre}}),
\]
using the protected secret key associated with $pk$. If $m\neq\bot$, it
signs $(\mathsf{sid},L_{\mathsf{pre}},m)$ using the signing component of $sk$
and appends
\[
  (\mathsf{Message},pk,m,\sigma)
\]
to $L$. Neither $A$ nor the unrevealed part of $sk$ is output. The ledger
retains the algorithm and its coins as the protected call record rather than
modelling an erasure.

If the owner is corrupt, the adversary knows $sk$ and may supply any pair of
bit strings $(m,\sigma)$ during the turn. The ledger appends
$(\mathsf{Message},pk,m,\sigma)$ verbatim without checking a signature, a key
relation, or any other protocol-specific predicate. If an
honest $A$ returns $\bot$, or if no message is supplied by the end of the
turn, the ledger appends $(\mathsf{NoMsg},pk)$. This timeout is immediately
visible to everyone and lets $\mathsf{WhosTurn}$ advance past a silent role.
After either record is appended, $pk$ is spent and is never scheduled again.

For an honest nonempty message, protected evaluation, signing, and append are
one atomic operation. The adversary sees the result only after it is
irrevocably in $L$. The ledger does not give an adversarial scheduler an
opportunity to inspect the result and then suppress it.

We assume that $\mathsf{KeyGen}$ may provide a public relation
$\mathsf{KeyMatch}(pk,sk)$. A protocol that uses disclosure of $sk$ as its
message can then check that the disclosed key belongs to the scheduled
public key. If the chosen key system has no such relation, the ledger does
not manufacture one: the protocol must tolerate the missing check or choose
different key syntax.

\subsection{How protocols use the interface}

The following examples explain the intended use without adding commands to
the ledger.

\paragraph{Ordinary computation role.}
From $\mathsf{crs}$ and the public list, all parties may derive the same
protocol object, including a ciphertext $C_{pk}$ under a scheduled public
key. This object need not itself be an entry of $L$. The honest role supplies
an algorithm $A$ which derives $C_{pk}$, decrypts its local state with $sk$,
examines $L$, and returns $sk$ if the role is activated and $\bot$ otherwise.
If it returns $sk$, everyone can check $\mathsf{KeyMatch}$, decrypt
$C_{pk}$, and continue the public computation. Thus publishing the key is the
role's only message and models complete post-speech exposure. If $A$ returns
$\bot$, the secret key and encrypted role state remain hidden forever.

For this interpretation, every secret used by an ordinary role must be
recoverable from $sk$ and public protocol data. A protocol that gives the
role additional hidden auxiliary state cannot claim that publishing only
$sk$ models leakage of its complete state.

\paragraph{Strong input role.}
An application party $P$ embeds its private input in $A$ and causes the
ledger to append only the prescribed value $A(sk_P,L)$ and its signature.
For example, $A$ may decrypt two activation tokens and return the one selected
by an input bit. If $P$ remains honest, neither $sk_P$, the unused token, nor
the input-dependent algorithm is revealed. If $P$ is corrupted, the
adversary receives its ordinary UC state together with the key and protected
call history retained by the ledger; correspondingly the simulator receives
the input leakage specified by the target functionality. Corruption after
the turn does not permit a replacement message because $pk_P$ is already
spent.

The present finite construction has public output. A later private-output
variant may add a dual protected-read operation which evaluates a function of
$(sk_P,L)$ and returns only its result privately to $P$; no such operation is
needed here.

\subsection{Random corruption and finite committees}

The random coins apply only to computation-party registrations. Let $C$ be
any size-$k$ computation committee selected from registration indices without
using the hidden coins $(c_i)_i$, and let
$X=|\{i\in C:c_i=1\}|$. Then $X\sim\mathsf{Binomial}(k,p)$. At the concrete
choice $p=1/4$, Hoeffding's inequality gives
\[
  \Pr[X\geq k/3]\leq \exp(-k/72).
\]
If a finite execution uses $q$ such committees, a union bound gives
\[
  \Pr[\neg\mathsf{Good}]\leq q\exp(-k/72),
\]
where $\mathsf{Good}$ is the event that every committee has at most
$t=\lfloor(k-1)/3\rfloor$ corrupt members. Hence
$k=\Theta(\lambda+\log q)$ makes the bad event negligible for every
polynomial-size finite execution. Application parties do not enter this
calculation.

\subsection{Interpretation and limitations}

The preliminary experiment should not be read as saying that a real adversary
naturally corrupts machines by tossing independent coins. The intended
real-world picture has a large population of apparently interchangeable
machines and an adversary that may choose any fixed $p$ fraction of that
population to compromise. Fresh computation clients, keys, and role states
are placed on machines at random, without revealing the placement before a
machine is corrupted or speaks. A key's logical role may be public even though
the physical machine carrying its secret state is not linkable from the
ledger. Random placement therefore turns the adversary's fixed physical set
into an approximately random corrupt subset of each computation committee.
The independent Bernoulli coins in the present model collapse this placement
argument into the cleaner statement that each computation PID is corrupt
independently.

This is a fixed-population, non-mobile abstraction. Once a computation PID is
sampled honest it cannot later be corrupted. An adversary that can repeatedly
replace the corrupted quarter of the physical population is strictly stronger
and is not modelled here. Named application parties are different on purpose:
an account owner or other designated input provider is not plausibly hidden by
random placement, so it retains the ordinary adaptive corruption interface.
The random computation corruptions are an implementation assumption to be
simulated away, not a corruption rule for the target application.

The statement that the ledger stores $sk$ and evaluates $A(sk,L)$ is likewise
an ideal-interface device, not a proposal that a blockchain hold users' secret
keys or execute their private code. In the intended implementation, a machine
stores its own key and locally computes the one value it places on the public
board. Moving that local step behind the ledger interface prevents an honest
party's secret key from becoming an output visible to the UC environment. For
a computation role, publishing $sk$ is an explicit way to model that speaking
exposes all of its encrypted role state; for a named application party, only
the prescribed value is published unless the party is actually corrupted.

This abstraction is intentionally stronger and cleaner than such a system.
It assumes honest key generation, independent placement, no selective abort
of an unfavourable registration, instantaneous common knowledge of $L$,
unpreventable honest appends, exact synchronous timeouts, protected
evaluation under a secret key, and a protocol-dependent turn function fixed
at initialization. It does not model consensus, denial of registration, network
partitions, grinding, compromise correlations, or the mechanism that places
a computation key on a real machine. These assumptions must not be mistaken
for properties already obtained by a concrete blockchain.

Finally, the computation PIDs and their random corruption coins exist only in
the real or ledger-hybrid execution. The target ideal functionality has the
usual application interface: named input and output parties, their prescribed
adaptive corruptions, and the intended outputs. It contains no committees,
randomly corrupted computation population, garbled circuits, or
$\mathsf{WhosTurn}$. A security proof must simulate all of that protocol
machinery from the target functionality's leakage.

The final paper may keep only this interface and interpretation in the main
text. Once the API is stable, an appendix can formulate
$\mathcal F_{\mathsf{Ledger}}$ in UC notation and make precise ITM creation,
private delivery and retention of $A$, synchronous timeouts, and corruption
bookkeeping. Those details should express the API above rather than determine
it.

\section{The Basic Non-Reactive Construction}
\label{sec:basic-construction}

We illustrate the central idea for a single evaluation of a public Boolean
circuit $F$. For now the output is public. The construction assumes an ideal
setup procedure capable of sampling hidden committees and anonymously
delivering their local state. Its purpose is to isolate how silence can
replace erasure, not yet to optimize assumptions or efficiency.

\subsection{Setup}

The setup procedure garbles $F$, obtaining a public garbled circuit
$\widetilde F$, two labels $(L^w_0,L^w_1)$ for every input wire $w$, and public
output-decoding information. For every input wire $w$ it also samples two
independent random activation tokens $(T^w_0,T^w_1)$.

For each pair $(w,b)$, setup selects a fresh hidden committee $C^w_b$. All
such committees are disjoint. It verifiably secret-shares $L^w_b$ among the
members of $C^w_b$ and anonymously delivers to each member its share, the wire
identifier, and a recognition value for $T^w_b$. A committee member learns
neither the bit $b$ represented by its label nor the identity of the sibling
committee. The input provider for wire $w$ receives the two tokens together
with the mapping $b\mapsto T^w_b$.

The garbled circuit, encrypted role packets, output-decoding information, and
ledger acceptance predicates are public. No online machine is given both
labels for an input wire.

\subsection{Input activation and evaluation}

To provide bit $x_w$, the input provider posts $T^w_{x_w}$ in the unique input
slot for wire $w$. The ledger authenticates the provider and accepts at most
one token in that slot. Every machine scans the accepted ledger. Precisely
the members of $C^w_{x_w}$ recognize the token and each posts its share as its
single lifetime message. Valid shares reconstruct $L^w_{x_w}$ on the ledger.
The sibling committee $C^w_{1-x_w}$ remains silent forever.

Once one input label has been reconstructed for every wire, anyone evaluates
$\widetilde F$ and decodes the public output. Public evaluation may itself be
performed by a fresh machine: its later exposure reveals only public data.

\subsection{Why post-speech exposure need not be fatal}

After a selected committee member speaks, its entire state reveals its label
share and activation information. The share belongs to a label that has
already been reconstructed publicly, so this leakage is harmless. Because
the two committees are disjoint, the exposed state contains no share of the
unselected label. That label remains protected by threshold privacy inside a
committee whose honest members never speak.

The same observation defeats the most immediate reset attack. Corrupting an
input provider after its post may reveal both activation tokens, but the
ledger will not accept a second token for the same input slot. Silent
committees react only to the uniquely accepted ledger value, not to arbitrary
strings that appear elsewhere.

This argument is only local. A complete proof must combine anonymous
recipient security, the hidden-assignment tail bound, robust threshold
reconstruction, garbling privacy, and simulation of the special input
providers. It must also rule out role collisions and setup grinding. The
blueprint is nevertheless enough to show why the model is not immediately
empty: a selected branch may become completely exposed while its mutually
exclusive sibling remains information-theoretically hidden by silence.

\paragraph{A gate-by-gate variant.}
For intuition, garbling can be replaced by a very expensive
information-theoretic construction. For every NAND gate and every pair of
input bits, a separate hidden committee holds shares of the corresponding
output token. Once the two accepted input tokens appear, only the matching
row committee reconstructs. This makes the role of silence especially
transparent, at the cost of a large setup and many committees.

\section{Toward Ongoing Computation}
\label{sec:ongoing}

The preceding construction is prepared for one circuit evaluation. An
ongoing computation must turn the opaque output labels of one garbling into
the input labels of a fresh garbling. Reactive garbling provides exactly such
links \cite{NielsenRanellucci2015ReactiveGarbling}, but the usual linker is
computed from both old output labels and both new input labels. A YAMSO
machine that computes and publishes this link would reveal all of these labels
when its state leaks.

Preparing all links in advance avoids the exposure problem only for a bounded
number of epochs. Asking the universal sampler for an independent fresh
garbling at every epoch does not solve the semantic problem either: the new
sample sees an opaque old label but has no correlated information telling it
whether that label represents zero or one.

We therefore propose a stronger object, provisionally called a
\emph{universal streamer}. One initial setup sample chooses a master key $K$
and publishes a protected program $PP_K$. Subsequent public evaluations
derive domain-separated label keys, branch packets, and hidden role
assignments for successive epochs. Schematically,
\[
  PP_K \leftarrow \mathsf{SetupStream}(1^\lambda,F),
  \qquad
  \mathsf{Prepare}(PP_K,i,D_i)
     \longrightarrow (\text{encrypted branch packets for epoch }i),
\]
where $D_i$ is an authenticated opaque state. The program prepares both
branch-local packets anonymously, but the accepted input token still causes
only one hidden committee to release its packet. Thus $PP_K$ supplies
cross-epoch correlation while silence continues to enforce one-of-two
disclosure.

There are close ingredients in the literature. Succinct garbling derives
adjacent garbling keys from one PRF seed inside a public obfuscated generator
\cite{LinPass2014SuccinctGarbling}. Cryptographic iterators publish an
obfuscated update program containing a hidden puncturable-PRF key and support
a very long sequence of compact state transitions
\cite{KoppulaLewkoWaters2014IOTM}. Streaming functional encryption formalizes
fresh inputs and evolving hidden state
\cite{GuanKorbSahai2022StreamingFE,BhushanKorbSahai2024StreamingFE,DattaGuanKorbSahai2024AdaptiveStreamingFE}.
Public obfuscation has also been
used to obtain adaptive security without erasures by ensuring that a speaking
party never possesses both label alternatives
\cite{CanettiGoldwasserPoburinnaya2014Adaptive2PC}.

None of these results directly provides the YAMSO interface. The missing
proof must combine authenticated canonical ledger paths, deterministic and
non-grindable role preparation, anonymous threshold delivery, fresh-machine
separation, and complete post-speech exposure. A useful development order is
therefore:
\begin{enumerate}
  \item prove the two-token, one-input construction formally;
  \item extend it to a finite sequence of pre-correlated transitions;
  \item formulate the universal-streamer functionality; and
  \item replace the finite table by a PRF-correlated public program.
\end{enumerate}

Reusable garbled circuits and garbled Turing machines remain relevant
comparison points
\cite{GoldwasserEtAl2012ReusableGarbling,GoldwasserEtAl2013TuringMachines}.
Their compact or reusable encodings do not
by themselves solve YAMSO: the standard encoder retains secret state whose
later exposure is outside the security guarantee. Their machinery may,
however, help realize the protected transition program once the desired
interface is stated precisely.


\bibliographystyle{splncs04}
\bibliography{bibliography}

\end{document}
